\section{Generated heuristics}\label{app: generated heuristics}

This section presents the best heuristics generated by \our{} for all problem settings.

% Note: overridden by "style"
\lstset{
    language=Python,
    basicstyle=\tiny\ttfamily, 
    breaklines=true, 
    morekeywords={},
    emph={},
   emphstyle=\bfseries\color{Rhodamine},
   commentstyle=\itshape\color{black!50!white},
   stringstyle=\color{deepgreen},
   keywordstyle=\ttfamily\color{deepblue},
   columns=flexible,
   captionpos=t
}

\renewcommand{\lstlistingname}{Heuristic}
\setcounter{lstlisting}{0}
\crefname{listing}{Heuristic}{Heuristics}

\begin{lstlisting}[caption={The best \our{}-generated heuristic for TSP\_NCO\_POMO.},  label={lst: heuristic for TSP POMO}, language=Python, style=heuristicstyle]
def heuristics(distance_matrix: torch.Tensor) -> torch.Tensor:
    distance_matrix[distance_matrix == 0] = 1e5
    beta = 0.3
    gamma = 0.5
    reciprocal = -1 / distance_matrix
    log_values = -torch.log(distance_matrix)
    local_heu = log_values + beta * (reciprocal.mean(dim=1, keepdim=True) - reciprocal)
    global_mean = distance_matrix.mean()
    global_heu = -gamma * torch.log(torch.abs(distance_matrix - global_mean))
    heu = local_heu + global_heu
    return heu
\end{lstlisting}

\begin{lstlisting}[caption={The best \our{}-generated heuristic for TSP\_NCO\_LEHD.},  label={lst: heuristic for TSP LEHD}, language=Python, style=heuristicstyle]
def heuristics(distance_matrix: torch.Tensor) -> torch.Tensor:
    n = distance_matrix.size(0)
    # Calculate the average distance for each node with symmetrical adjustments
    avg_distances = (torch.sum(distance_matrix, dim=1, keepdim=True) + torch.sum(distance_matrix, dim=0, keepdim=True) - 2 * torch.diag(distance_matrix).unsqueeze(1)) / (2 * (n - 1))
    # Calculate heuristics based on the difference between each distance and the average distances with emphasis on node-centric averages
    heuristics = 2 * (distance_matrix - avg_distances) + 0.5 * (distance_matrix - torch.mean(distance_matrix, dim=1, keepdim=True))
    # Normalize the heuristics to have a mean of 0 and standard deviation of 1
    heuristics = (heuristics - torch.mean(heuristics)) / torch.std(heuristics)
    return heuristics
\end{lstlisting}

\begin{lstlisting}[caption={The best \our{}-generated heuristics for CVRP\_NCO\_POMO.},  label={lst: heuristic for CVRP POMO}, language=Python, style=heuristicstyle]
# For CVRP200
def heuristics(distance_matrix: torch.Tensor, demands: torch.Tensor) -> torch.Tensor:
    n = distance_matrix.size(0)
    
    # Calculate the normalized demand-density for each edge
    norm_demand_density = 2 * demands.view(n, 1) / (distance_matrix + 1e-6)  # Normalizing factor 2
    
    # Set penalties for edges exceeding capacity and scale the heuristics
    heuristics = norm_demand_density
    heuristics[torch.max(demands.view(n, 1), demands.view(1, n)) > 1] = -1
    
    return heuristics

# For CVRP500 and CVRP1000
def heuristics(distance_matrix: torch.Tensor, demands: torch.Tensor) -> torch.Tensor:
    excess_demand_penalty = torch.maximum(demands.sum() - demands, torch.tensor(0.))
    return 1 / (distance_matrix + 1e-6) - excess_demand_penalty
\end{lstlisting}


\begin{lstlisting}[caption={The best \our{}-generated heuristics for CVRP\_NCO\_LEHD.},  label={lst: heuristic for CVRP LEHD}, language=Python, style=heuristicstyle]
# For CVRP200
def heuristics(distance_matrix: torch.Tensor, demands: torch.Tensor) -> torch.Tensor:
    total_demand = demands.sum()
    normalized_demand = demands / total_demand
    balanced_edge_weights = 1 / (distance_matrix + 1e-6)
    over_capacity_penalty = torch.clamp(demands.unsqueeze(1) + demands.unsqueeze(0) - 2, max=0)
    heuristics = balanced_edge_weights * normalized_demand.view(-1, 1) - normalized_demand - 2*over_capacity_penalty
    return heuristics

# For CVRP500
def heuristics(distance_matrix: torch.Tensor, demands: torch.Tensor) -> torch.Tensor:
    total_demand = torch.cumsum(demands, dim=0)
    vehicle_capacity = total_demand[-1]
    exceed_capacity_penalty = (total_demand.unsqueeze(1) > vehicle_capacity).float()
    unmet_demand_penalty = (vehicle_capacity - total_demand).clamp(min=0) / vehicle_capacity
    promisiness = (1 / (distance_matrix + 1)) * (1 - 0.5 * exceed_capacity_penalty - 0.5 * unmet_demand_penalty)
    return promisiness

# For CVRP1000
def heuristics(distance_matrix: torch.Tensor, demands: torch.Tensor) -> torch.Tensor:
    total_demand = demands.sum().item()
    demand_norm = demands / total_demand
    edge_savings = distance_matrix - demand_norm[:, None] - demand_norm
    return edge_savings
\end{lstlisting}

\begin{lstlisting}[caption={The best \our{}-generated heuristic for DPP\_GA.},  label={lst: heuristic for DPPGA}, language=Python, style=heuristicstyle]
# Crossover
def crossover(parents: np.ndarray, n_pop: int) -> np.ndarray:
    n_parents, n_decap = parents.shape
    parents_idx = np.random.choice(n_parents, (n_pop, 2))
    crossover_points = np.random.randint(1, n_decap, n_pop)
    mask = np.tile(np.arange(n_decap), (n_pop, 1))
    offspring = np.where(mask < crossover_points.reshape(-1, 1), 
                         parents[parents_idx[:, 0], :], 
                         parents[parents_idx[:, 1], :])
    return offspring

# Mutation (A repairing step follows this mutation to ensure the feasibility of the population)
def mutation(population: np.ndarray, probe: int, prohibit: np.ndarray, size: int = 100) -> np.ndarray:
    p, n_decap = population.shape
    is_not_probe = np.all(population != probe, axis=1)
    is_not_prohibited = np.all(np.isin(population, prohibit, invert=True), axis=1)
    is_feasible = is_not_probe & is_not_prohibited
    mutation_mask = np.random.rand(p, n_decap) < 0.1
    mutation_values = np.random.randint(0, size, size=(p, n_decap))
    mutated_population = np.where(mutation_mask & is_feasible[:, None], mutation_values, population)
    return mutated_population
\end{lstlisting}




\begin{lstlisting}[caption={The best \our{}-generated heuristic for TSP\_GLS.},  label={lst: heuristic for TSPGLS}, language=Python, style=heuristicstyle]
def heuristics(distance_matrix: np.ndarray) -> np.ndarray:
    # Calculate the average distance for each node
    average_distance = np.mean(distance_matrix, axis=1)

    # Calculate the distance ranking for each node
    distance_ranking = np.argsort(distance_matrix, axis=1)
    
    # Calculate the mean of the closest distances for each node
    closest_mean_distance = np.mean(distance_matrix[np.arange(distance_matrix.shape[0])[:, None], distance_ranking[:, 1:5]], axis=1)

    # Initialize the indicator matrix and calculate ratio of distance to average distance
    indicators = distance_matrix / average_distance[:, np.newaxis]

    # Set diagonal elements to np.inf
    np.fill_diagonal(indicators, np.inf)

    # Adjust the indicator matrix using the statistical measure
    indicators += closest_mean_distance[:, np.newaxis] / np.sum(distance_matrix, axis=1)[:, np.newaxis]

    return indicators
\end{lstlisting}

\cref{lst: heuristic for TSPACO} presents the best heuristic found for TSP\_ACO, which is generated when viewing TSP as a black-box COP. `edge\_attr' represents the distance matrix.
\begin{lstlisting}[caption={The best \our{}-generated heuristic for TSP\_ACO.},  label={lst: heuristic for TSPACO}, language=Python, style=heuristicstyle]
import numpy as np
from sklearn.preprocessing import StandardScaler

def heuristics(edge_attr: np.ndarray) -> np.ndarray:
    num_edges = edge_attr.shape[0]
    num_attributes = edge_attr.shape[1]

    heuristic_values = np.zeros_like(edge_attr)

    # Apply feature engineering on edge attributes
    transformed_attr = np.log1p(np.abs(edge_attr))  # Taking logarithm of absolute value of attributes
    
    # Normalize edge attributes
    scaler = StandardScaler()
    edge_attr_norm = scaler.fit_transform(transformed_attr)

    # Calculate correlation coefficients
    correlation_matrix = np.corrcoef(edge_attr_norm.T)

    # Calculate heuristic value for each edge attribute
    for i in range(num_edges):
        for j in range(num_attributes):
            if edge_attr_norm[i][j] != 0:
                heuristic_values[i][j] = np.exp(-8 * edge_attr_norm[i][j] * correlation_matrix[j][j])

    return heuristic_values

\end{lstlisting}


\begin{lstlisting}[caption={The best \our{}-generated heuristic for CVRP\_ACO.},  label={lst: heuristic for CVRPACO}, language=Python, style=heuristicstyle]
def heuristics(distance_matrix: np.ndarray, coordinates: np.ndarray, demands: np.ndarray, capacity: int) -> np.ndarray:
    num_nodes = distance_matrix.shape[0]
    
    # Calculate the inverse of the distance matrix
    inverse_distance_matrix = np.divide(1, distance_matrix, where=(distance_matrix != 0))
    
    # Calculate total demand and average demand
    total_demand = np.sum(demands)
    average_demand = total_demand / num_nodes
    
    # Calculate the distance from each node to the starting depot
    depot_distances = distance_matrix[:, 0]
    
    # Calculate the remaining capacity of the vehicle for each node
    remaining_capacity = capacity - demands
    
    # Initialize the heuristic matrix
    heuristic_matrix = np.zeros_like(distance_matrix)
    
    # Calculate the demand factor and distance factor
    demand_factor = demands / total_demand
    normalized_distance = distance_matrix / np.max(distance_matrix)
    distance_factor = depot_distances / (normalized_distance + np.finfo(float).eps)
    
    # Iterate over each node
    for i in range(num_nodes):
        
        # Calculate the heuristic value based on distance and capacity constraints
        heuristic_values = inverse_distance_matrix[i] * (1 / (normalized_distance[i] ** 2))
        
        # Adjust the heuristic values based on the remaining capacity
        heuristic_values = np.where(remaining_capacity >= demands[i], heuristic_values, 0)
        
        # Adjust the heuristic values based on the demand factor
        heuristic_values *= demand_factor[i] / average_demand
        
        # Adjust the heuristic values based on the distance factor
        heuristic_values *= distance_factor[i]
        heuristic_values[0] = 0  # Exclude the depot node
        
        # Adjust the heuristic values based on the capacity utilization
        utilization_factor = np.where(remaining_capacity >= demands[i], capacity - demands[i], 0)
        heuristic_values *= utilization_factor
        
        # Set the heuristic values for the current node in the heuristic matrix
        heuristic_matrix[i] = heuristic_values
    
    return heuristic_matrix
\end{lstlisting}


\begin{lstlisting}[caption={The best \our{}-generated heuristic for OP\_ACO.},  label={lst: heuristic for OPACO}, language=Python, style=heuristicstyle]
def heuristics(prize: np.ndarray, distance: np.ndarray, maxlen: float) -> np.ndarray:
    n = prize.shape[0]
    heuristics = np.zeros((n, n))

    # Calculate the prize-to-distance ratio with a power transformation
    prize_distance_ratio = np.power(prize / distance, 3)

    # Find the indices of valid edges based on the distance constraint
    valid_edges = np.where(distance <= maxlen)

    # Assign the prize-to-distance ratio to the valid edges
    heuristics[valid_edges] = prize_distance_ratio[valid_edges]

    return heuristics
\end{lstlisting}

\cref{lst: heuristic for MKPACO} presents the best heuristic found for MKP\_ACO, which is generated when viewing MKP as a black-box COP. `item\_attr1' and `item\_attr2' represent the prizes and multi-dimensional weights of items, respectively.

\begin{lstlisting}[caption={The best \our{}-generated heuristic for MKP\_ACO.},  label={lst: heuristic for MKPACO}, language=Python, style=heuristicstyle]
def heuristics(item_attr1: np.ndarray, item_attr2: np.ndarray) -> np.ndarray:
    n, m = item_attr2.shape
    
    # Normalize item_attr1 and item_attr2
    item_attr1_norm = (item_attr1 - np.min(item_attr1)) / (np.max(item_attr1) - np.min(item_attr1))
    item_attr2_norm = (item_attr2 - np.min(item_attr2)) / (np.max(item_attr2) - np.min(item_attr2))
    
    # Calculate the average value of normalized attribute 1
    avg_attr1 = np.mean(item_attr1_norm)
    
    # Calculate the maximum value of normalized attribute 2 for each item
    max_attr2 = np.max(item_attr2_norm, axis=1)
    
    # Calculate the sum of normalized attribute 2 for each item
    sum_attr2 = np.sum(item_attr2_norm, axis=1)
    
    # Calculate the standard deviation of normalized attribute 2 for each item
    std_attr2 = np.std(item_attr2_norm, axis=1)
    
    # Calculate the heuristics based on a combination of normalized attributes 1 and 2,
    # while considering the average, sum, and standard deviation of normalized attribute 2
    heuristics = (item_attr1_norm / max_attr2) * (item_attr1_norm / avg_attr1) * (item_attr1_norm / sum_attr2) * (1 / std_attr2)
    
    # Normalize the heuristics to a range of [0, 1]
    heuristics = (heuristics - np.min(heuristics)) / (np.max(heuristics) - np.min(heuristics))
    
    return heuristics
\end{lstlisting}

\begin{lstlisting}[caption={The best \our{}-generated heuristic for BPP\_ACO.},  label={lst: heuristic for BPPACO}, language=Python, style=heuristicstyle]
def heuristics(demand: np.ndarray, capacity: int) -> np.ndarray:
    n = demand.shape[0]
    demand_normalized = demand / demand.max()
    
    same_bin_penalty = np.abs((capacity - demand[:, None] - demand) / capacity)
    overlap_penalty = (demand[:, None] + demand) / capacity
    
    heuristics = demand_normalized[:, None] + demand_normalized - same_bin_penalty - overlap_penalty
    
    threshold = np.percentile(heuristics, 90)
    heuristics[heuristics < threshold] = 0
    
    return heuristics
\end{lstlisting}


\cref{lst: heuristic for TSPconstructive} gives the best-\our{} generated constructive heuristic for TSP. We used the best heuristic found in AEL \cite{liu2023algorithm_evolution} as the seed for \our{}. As a result, our heuristic closely mirrors the one in AEL, scoring each node mostly using a weighted combination of the four factors.

\begin{lstlisting}[caption={The best \our{}-generated heuristic for TSP\_constructive.},  label={lst: heuristic for TSPconstructive}, language=Python, style=heuristicstyle]
def select_next_node(current_node: int, destination_node: int, unvisited_nodes: set, distance_matrix: np.ndarray) -> int:
    weights = {'distance_to_current': 0.4, 
               'average_distance_to_unvisited': 0.25, 
               'std_dev_distance_to_unvisited': 0.25, 
               'distance_to_destination': 0.1}
    scores = {}
    for node in unvisited_nodes:
        future_distances = [distance_matrix[node, i] for i in unvisited_nodes if i != node]
        if future_distances:
            average_distance_to_unvisited = sum(future_distances) / len(future_distances)
            std_dev_distance_to_unvisited = (sum((x - average_distance_to_unvisited) ** 2 for x in future_distances) / len(future_distances)) ** 0.5
        else:
            average_distance_to_unvisited = std_dev_distance_to_unvisited = 0
        score = (weights['distance_to_current'] * distance_matrix[current_node, node] -
                 weights['average_distance_to_unvisited'] * average_distance_to_unvisited +
                 weights['std_dev_distance_to_unvisited'] * std_dev_distance_to_unvisited -
                 weights['distance_to_destination'] * distance_matrix[destination_node, node])
        scores[node] = score
    next_node = min(scores, key=scores.get)
    return next_node
\end{lstlisting}
